\section*{Abstract}


This thesis explores the possibility of developing a multi-user neurofeedback system designed for classrooms, using low-cost, single-electrode Brain-Computer Interface devices (BCI). We opted to use the NeuroSky MindWave Mobile 2 headsets for this study. The system will aim to collect and process attention data from multiple students in real time, giving educators an insightful overview of student focus and engagement throughout different segments of a lecture.


The architecture of the system is built around a network of BCI headsets transmitting data through a Lab Streaming Layer (LSL) protocol, to a computer that will be operated by the teacher in the front, something that most personnel already have at hand. A dedicated Python application will also be implemented, with the goal of visualizing each student’s attention level dynamically, enabling the instructor to identify not only patterns of engagement, but also of disengagement. The system also allows lectures to be segmented into predefined sections, correlating attention levels with specific parts of the material. This system helps with post-lecture analysis to determine which topics or segments sustained the most (or least) student interest, supporting educators to do data-driven adjustments in their teaching strategies.


Contrary to most papers on this theme, the focus of this thesis is not on analyzing learning outcomes or comparing teaching methodologies, but rather on evaluating the technical and practical feasibility of such a system in a classroom context and how it can influence learning as a whole. This work investigates aspects such as data synchronization, signal quality, scalability, and user experience, giving us a foundation for future research in neurofeedback-based education utilizing this system. Ultimately, this project aims to contribute toward the development of intelligent educational tools that leverage neurotechnology to improve teaching effectiveness and engagement of students.